% ============================================================
%  DIAPOSITIVAS – Algoritmos Genéticos para Termas Solares
%  Formato: Beamer, tema Madrid limpio, sin cajas raras
% ============================================================
\documentclass[10pt]{beamer}

\usetheme{Madrid}
\usecolortheme{default}

% Paleta de colores personalizada (azul oscuro profesional)
\definecolor{azulunsa}{RGB}{0,56,101}
\definecolor{azulclaro}{RGB}{0,100,180}
\definecolor{grisclaro}{RGB}{240,240,240}

\setbeamercolor{palette primary}{bg=azulunsa,fg=white}
\setbeamercolor{palette secondary}{bg=azulclaro,fg=white}
\setbeamercolor{palette tertiary}{bg=azulunsa,fg=white}
\setbeamercolor{palette quaternary}{bg=azulclaro,fg=white}
\setbeamercolor{structure}{fg=azulunsa}
\setbeamercolor{frametitle}{bg=azulunsa,fg=white}
\setbeamercolor{title}{bg=azulunsa,fg=white}
\setbeamercolor{block title}{bg=azulunsa,fg=white}
\setbeamercolor{block body}{bg=grisclaro,fg=black}
\setbeamercolor{itemize item}{fg=azulclaro}
\setbeamercolor{itemize subitem}{fg=azulclaro}

\setbeamertemplate{navigation symbols}{}
\setbeamertemplate{footline}[frame number]

% Paquetes
\usepackage[utf8]{inputenc}
\usepackage[T1]{fontenc}
\usepackage[spanish]{babel}
\usepackage{amsmath}
\usepackage{amssymb}
\usepackage{graphicx}
\usepackage{booktabs}
\usepackage{listings}
\usepackage{xcolor}

\lstset{
  language=Python,
  basicstyle=\ttfamily\tiny,
  keywordstyle=\color{azulclaro}\bfseries,
  stringstyle=\color{red!70!black},
  commentstyle=\color{green!50!black}\itshape,
  frame=single,
  rulecolor=\color{azulunsa},
  breaklines=true,
  showstringspaces=false,
  numbers=left,
  numberstyle=\tiny\color{gray},
  numbersep=5pt
}

% ============================================================
%  PORTADA
% ============================================================
\title[AG Termas Solares]{Algoritmos Genéticos para Determinar el Ángulo de\\
       Inclinación Óptimo de Termas Solares en el Perú}
\subtitle{Réplica, implementación desde cero y mejora memética}
\author{Edilson Bonet Mamani Yucra}
\institute[UNSA]{Universidad Nacional de San Agustín de Arequipa\\
                 Escuela Profesional de Ciencia de la Computación\\[0.3em]
                 \textit{Docente: Rolando Cardenas Talavera}}
\date{Arequipa, \today}

\begin{document}

% ── Portada ──────────────────────────────────────────────────
\begin{frame}
  \titlepage
\end{frame}

% ── Índice ───────────────────────────────────────────────────
\begin{frame}{Contenido}
  \tableofcontents
\end{frame}

% ============================================================
\section{Introducción y Motivación}
% ============================================================
\begin{frame}{Motivación del Problema}
  \begin{itemize}
    \item El Perú tiene uno de los mayores potenciales de radiación solar del planeta.
    \item Las termas solares de tubos al vacío son sistemas de bajo costo para calentar agua.
    \item Su eficiencia depende del \textbf{ángulo de inclinación $\beta$} del colector.
    \item En la práctica, $\beta$ se fija empíricamente, sin considerar la geografía local.
  \end{itemize}

  \vspace{0.4cm}

  \begin{block}{Objetivo}
    Replicar íntegramente, \textbf{desde cero}, la metodología de una tesis de\\
    Ingeniería Mecatrónica que aplica Algoritmos Genéticos para determinar\\
    automáticamente el ángulo óptimo $\beta^*$ en 6 ciudades del Perú.\\[0.2em]
    \textbf{Sin} librerías de IA (ni DEAP, ni scipy). Solo Python estándar.
  \end{block}
\end{frame}

% ============================================================
\section{Datos del Problema (Dataset)}
% ============================================================
\begin{frame}{Dataset: 6 Ciudades del Perú}
  \small
  \begin{table}
    \centering
    \caption{Coordenadas geográficas y parámetros de los tubos al vacío}
    \begin{tabular}{lrrr}
      \toprule
      \textbf{Ciudad} & \textbf{Latitud (°)} & \textbf{Longitud (°)} & \textbf{Altitud (km)} \\
      \midrule
      Trujillo  & $-8.1117$  & $-79.0286$ & 0.034 \\
      Piura     & $-5.1971$  & $-80.6266$ & 0.055 \\
      Iquitos   & $-3.7493$  & $-73.2444$ & 0.106 \\
      Tacna     & $-18.0137$ & $-70.2510$ & 0.552 \\
      Arequipa  & $-16.4154$ & $-71.5218$ & 2.335 \\
      Puno      & $-15.8405$ & $-70.0279$ & 3.827 \\
      \bottomrule
    \end{tabular}
  \end{table}

  \vspace{0.2cm}
  \textbf{Parámetros fijos del tubo:}
  $D_1 = 47\,\text{mm}$, $D_2 = 58\,\text{mm}$,
  $L_{tubo} = 1.80\,\text{m}$, $B = 80\,\text{mm}$
\end{frame}

% ============================================================
\section{Modelo Matemático (Función Objetivo)}
% ============================================================
\begin{frame}{Modelo de Geometría Solar — Función Objetivo}
  La función objetivo es la \textbf{energía anual} recolectable por unidad de longitud:

  \begin{align*}
    H_{anual} &= \sum_{n=1}^{365}
      \int_{-\omega_{ss}}^{\omega_{ss}} P_{total}(t)\, dt
      \quad \text{[ecuación 34 de la tesis]}\\[4pt]
    P_{total} &= D_1 L \Bigl(G_{bn}\cos\theta_t\, f(\Omega)
                           + \pi\, G_{d\beta}\, F\Bigr)
      \quad \text{[ec. 32]}\\[4pt]
    \tau_b &= r_0 a_0 + r_1 a_1 \exp\!\Bigl(\tfrac{-r_k k}{\cos\theta_Z}\Bigr)
      \quad \text{(Hottel, clima tropical)}
  \end{align*}

  \vspace{0.2cm}
  \begin{itemize}
    \item $\omega_{ss}$: ángulo horario de atardecer (superficie inclinada).
    \item $f(\Omega)$: función de aceptación angular del tubo.
    \item $F$: factor de forma difuso. Ambos dependen de $D_1, D_2, B$.
    \item Integración: \textbf{Regla de Simpson compuesta}, 241 puntos/día.
  \end{itemize}
\end{frame}

\begin{frame}[fragile]{Implementación: Cálculo de Energía Anual}
  Archivo \texttt{modelo\_solar.py} — función objetivo que maximiza el AG:

  \begin{lstlisting}
def energia_anual_MJ(phi_deg, altitud_km, beta_deg, B_mm,
                     paso_dias=1, n_pasos=241):
    """Energia anual [MJ/m] recolectada por el tubo."""
    total = 0.0
    dias = range(1, 366, paso_dias)
    for n in dias:
        total += energia_diaria_MJ(n, phi_deg, altitud_km,
                                   beta_deg, B_mm,
                                   n_pasos=n_pasos)
    if paso_dias > 1:
        total *= 365.0 / len(dias)
    return total
  \end{lstlisting}

  \begin{itemize}
    \item Itera los 365 días del año.
    \item Integración numérica de Simpson por día.
    \item Acepta \texttt{paso\_dias} para acelerar el cómputo en el AG.
  \end{itemize}
\end{frame}

% ============================================================
\section{Algoritmo Genético Original (Réplica)}
% ============================================================
\begin{frame}{AG Original — Parámetros Replicados}
  \small
  \begin{table}
    \centering
    \caption{Parámetros del Algoritmo Genético original (réplica de la tesis)}
    \begin{tabular}{ll}
      \toprule
      \textbf{Parámetro} & \textbf{Valor / Método} \\
      \midrule
      Representación      & Valor real (ángulo $\beta$ en grados)           \\
      Función de aptitud  & $H_{anual}(\beta)$ (a maximizar)                 \\
      Tamaño de población & 50 individuos                                    \\
      Generaciones        & 15                                               \\
      Selección           & Torneo ($k = 3$)                                 \\
      Cruzamiento         & Blend Crossover ($\alpha=0.5$,\ $p_c=0.7$)      \\
      Mutación            & Gaussiana ($\sigma=10\%$ del rango,\ $p_m=0.1$) \\
      Reemplazo           & Generacional puro (sin elitismo)                 \\
      \bottomrule
    \end{tabular}
  \end{table}
\end{frame}

\begin{frame}[fragile]{AG Original — Selección por Torneo}
  Archivo \texttt{ga\_original.py}:

  \begin{lstlisting}
def seleccion_torneo(rng, poblacion, fitnesses, k_torneo=3):
    seleccionados = []
    n = len(poblacion)
    for _ in range(n):
        # Elegir k participantes al azar
        participantes = [rng.randrange(n) for _ in range(k_torneo)]
        # El que mayor fitness tenga gana
        ganador = max(participantes, key=lambda i: fitnesses[i])
        seleccionados.append(list(poblacion[ganador]))
    return seleccionados

def mutacion_gaussiana(rng, individuo, lo, hi, sigma, prob_mut=0.1):
    if rng.random() < prob_mut:
        individuo[0] += rng.gauss(0, sigma)
        # Mantener dentro de los limites de busqueda
        individuo[0] = max(lo, min(hi, individuo[0]))
  \end{lstlisting}
\end{frame}

% ============================================================
\section{Mejora Propuesta: Algoritmo Memético}
% ============================================================
\begin{frame}{Motivación de la Mejora}
  \begin{itemize}
    \item La propia tesis valida sus resultados comparándolos contra la\\
          \textbf{Búsqueda de la Sección Áurea} (GSS), un método unidimensional exacto.
    \item La función objetivo es \textbf{unimodal}: tiene un único máximo global.
    \item Esto abre la puerta a integrar la búsqueda local \emph{dentro} del AG.
  \end{itemize}

  \vspace{0.4cm}

  \begin{block}{Propuesta: Algoritmo Memético (Moscato, 1989)}
    Combinar exploración evolutiva global con explotación local exacta:\\[0.3em]
    \begin{enumerate}
      \item \textbf{Elitismo}: los 2 mejores individuos siempre sobreviven.
      \item \textbf{Mutación gaussiana adaptativa}: $\sigma$ decrece linealmente\\
            (de exploración amplia $\to$ explotación fina).
      \item \textbf{Refinamiento memético}: Búsqueda de la Sección Áurea\\
            cada 5 generaciones sobre el mejor individuo (aprendizaje Lamarckiano).
    \end{enumerate}
  \end{block}
\end{frame}

\begin{frame}[fragile]{Mejora — Búsqueda de la Sección Áurea (GSS)}
  Archivo \texttt{ga\_mejorado.py}:

  \begin{lstlisting}
def busqueda_seccion_aurea(fitness_fn, lo, hi, tol=1e-4, max_iter=60):
    """Maximiza fitness_fn en [lo, hi] usando la razon aurea."""
    gr = (5**0.5 - 1) / 2.0   # ~0.618
    a, b = lo, hi
    c = b - gr * (b - a)
    d = a + gr * (b - a)
    fc, fd = fitness_fn(c), fitness_fn(d)
    for _ in range(max_iter):
        if abs(b - a) < tol:
            break
        if fc > fd:
            b, d, fd = d, c, fc
            c = b - gr * (b - a);  fc = fitness_fn(c)
        else:
            a, c, fc = c, d, fd
            d = a + gr * (b - a);  fd = fitness_fn(d)
    return (a + b) / 2.0, fitness_fn((a + b) / 2.0)
  \end{lstlisting}

  Se aplica en un entorno de $\pm 8\%$ del rango alrededor del mejor individuo.
\end{frame}

% ============================================================
\section{Resultados}
% ============================================================
\begin{frame}{Resultados Comparativos — 6 Ciudades}
  \small
  \begin{table}
    \centering
    \caption{Ángulo óptimo $\beta^*$ y energía anual ($B=80$mm) vs. tesis original}
    \begin{tabular}{lrrrrr}
      \toprule
      \textbf{Ciudad} &
      \textbf{$\beta^*$ Tesis} &
      \textbf{$\beta^*$ AG Orig.} &
      \textbf{$\beta^*$ AG Mej.} &
      \textbf{$E$ Tesis} &
      \textbf{$E$ Réplica} \\
       & (°) & (°) & (°) & (MJ/m) & (MJ/m) \\
      \midrule
      Trujillo & 5.62  & 4.95  & 4.97  & 718.82 & 689.72 \\
      Piura    & 3.99  & 3.25  & 3.25  & 725.11 & 695.77 \\
      Iquitos  & 2.87  & 2.47  & 2.48  & 730.23 & 700.33 \\
      Tacna    & 13.64 & 11.75 & 11.77 & 719.45 & 685.19 \\
      Arequipa & 12.66 & 11.30 & 11.36 & 793.56 & 746.82 \\
      Puno     & 12.22 & 11.05 & 10.98 & 804.70 & 755.28 \\
      \bottomrule
    \end{tabular}
  \end{table}

  \vspace{0.2cm}
  \footnotesize
  Diferencia promedio en $\beta^*$: $\approx 1.1°$ \quad|\quad
  Energía réplica $\approx 95\%$ de la tesis (diferencia sistemática de calibración).
\end{frame}

\begin{frame}{Convergencia: AG Original vs. AG Mejorado (Trujillo)}
  \begin{figure}
    \centering
    \includegraphics[width=0.90\textwidth]{conver.png}
    \caption{Evolución del mejor y peor individuo por generación — Trujillo.
             El AG Mejorado concentra toda la población cerca del óptimo desde la generación 4.}
  \end{figure}
\end{frame}

\begin{frame}{Análisis de la Convergencia}
  \begin{itemize}
    \item Ambos AG encuentran prácticamente el \textbf{mismo valor óptimo final}.
    \item La diferencia es la \textbf{robustez de la población}:
  \end{itemize}

  \vspace{0.3cm}

  \begin{table}
    \centering
    \small
    \begin{tabular}{lrr}
      \toprule
      \textbf{Métrica (Trujillo, Gen. 15)} & \textbf{AG Original} & \textbf{AG Mejorado} \\
      \midrule
      Fitness mejor individuo (MJ/m) & 690.48 & 690.48 \\
      Fitness peor individuo (MJ/m)  & 686.72 & 690.45 \\
      Dispersión (mejor $-$ peor)    & 3.76   & 0.03   \\
      \bottomrule
    \end{tabular}
  \end{table}

  \vspace{0.3cm}
  \begin{itemize}
    \item El AG Mejorado reduce la dispersión en más del \textbf{99\%}.
    \item El barrido exhaustivo (fuerza bruta, paso $0.05°$) confirma el óptimo\\
          con diferencia $< 0.05°$ para las 6 ciudades.
    \item Iquitos $<$ Piura $<$ Trujillo $<$ Puno $\approx$ Arequipa $<$ Tacna:\\
          \textbf{el orden físico (latitud $\to$ $\beta^*$) se preserva perfectamente.}
  \end{itemize}
\end{frame}

% ============================================================
\section{Conclusiones}
% ============================================================
\begin{frame}{Conclusiones}
  \begin{enumerate}
    \item Se replicó íntegramente desde cero el modelo de geometría solar y el
          AG de la tesis, \textbf{sin librerías externas de IA}.
    \item El ángulo óptimo $\beta^*$ concuerda con la tesis (diferencia
          promedio $\approx 1.1°$) y conserva el mismo \textbf{orden relativo entre ciudades}.
    \item La diferencia en energía absoluta ($\approx 5\%$) es sistemática y
          atribuible a supuestos de calibración no declarados por el autor,
          \textbf{no a un error conceptual}.
    \item La mejora memética no altera el valor del óptimo (problema
          unidimensional simple) pero reduce la dispersión de la población en
          más del \textbf{99\%}, siendo especialmente relevante en problemas
          de mayor dimensionalidad.
    \item Se demuestra la viabilidad de aplicar AGs implementados desde cero
          a un problema real de ingeniería de energías renovables.
  \end{enumerate}
\end{frame}

\begin{frame}
  \centering
  \vspace{1.5cm}
  {\Large\bfseries Gracias por su atención}\\[1.5cm]
  {\normalsize Edilson Bonet Mamani Yucra}\\[0.3em]
  {\small Universidad Nacional de San Agustín de Arequipa}\\[0.3em]
  {\small Curso: Inteligencia Artificial}
\end{frame}

\end{document}
